
Puji syukur penulis panjatkan kehadirat Allah SWT, karena atas rahmat, hidayah, dan karunia-Nya, penulis dapat menyelesaikan penyusunan buku Proyek Akhir yang berjudul “Sistem Manajemen Praktik Kerja Lapangan (PKL) SMK PGRI Telagasari” dengan baik. Buku Proyek Akhir ini disusun sebagai salah satu syarat akademis untuk memperoleh gelar Ahli Madya (A.Md.Kom.) pada Program Studi Diploma Tiga Teknik Informatika, Departemen Teknik Informatika dan Komputer di Politeknik Elektronika Negeri Surabaya.

Penulis menyadari bahwa keberhasilan pelaksanaan penelitian hingga penyusunan buku Proyek Akhir ini tidak lepas dari bimbingan, arahan, motivasi, serta dukungan dari berbagai pihak. Oleh karena itu, pada kesempatan ini penulis menyampaikan apresiasi dan ucapan terima kasih yang sebesar-besarnya kepada:

\begin{enumerate}[label=\arabic*.]
\item Bapak Mohammad Robihul Mufid, M.Tr.Kom., selaku Dosen Pembimbing 1, serta Bapak Yanuar Risah Prayogi, S.Kom., M.Kom., selaku Dosen Pembimbing 2, yang telah meluangkan waktu, memberikan bimbingan teknis, serta arahan yang sangat berharga selama proses perancangan, pengembangan sistem InternSync, hingga penyusunan laporan ini selesai.
\item Bapak Direktur Politeknik Elektronika Negeri Surabaya beserta jajaran manajemen Departemen Teknik Informatika dan Komputer atas fasilitas dan iklim akademik yang baik selama penulis menempuh studi.
\item Pihak manajemen, koordinator, serta guru pembimbing SMK PGRI Telagasari yang telah memberikan izin penelitian, menyediakan data operasional, serta membantu memvalidasi kebutuhan sistem secara mendalam.
\item Ibu Ida Farida dan Ibu Odah selaku orang tua dan nenek tercinta, serta seluruh keluarga besar yang tiada hentinya memberikan dukungan moral, materiil, dan doa yang tulus demi kelancaran studi penulis.
\item Rekan-rekan mahasiswa Program Studi Teknik Informatika PENS 2023 serta seluruh pihak yang tidak dapat disebutkan satu per satu, atas bantuan, diskusi teknis, dan solidaritasnya selama penyelesaian proyek akhir ini.
\end{enumerate}

Penulis menyadari bahwa sistem InternSync maupun penulisan buku Proyek Akhir ini masih jauh dari kesempurnaan dan memiliki keterbatasan. Oleh karena itu, penulis mengharapkan kritik dan saran yang membangun dari pembaca demi pengembangan perangkat lunak ini di masa mendatang. Akhir kata, semoga Proyek Akhir ini dapat memberikan manfaat nyata bagi dunia pendidikan dan menjadi referensi ilmiah bagi perkembangan ilmu teknologi informasi.
